\par\Needspace{10\baselineskip}
\originalchapter{官方作业一}
本章翻译18.950 Homework 1的三个原题, 共20分. 原PDF没有标注截止日. 以下解答由AI独立编写并复审, 不是官方答案; 平面曲率采用第1章的有向约定.

\begin{exercise}\label{ps:1-1}
原题PS1.1, 5分. 正则曲线的曲率在平面 $\R^2$ 的下列变换下怎样变化: (a) 平移; (b) 旋转; (c) 反射; (d) 伸缩 $x\mapsto rx$?
\end{exercise}
\begin{solution}\label{pssol:1-1}
设 $c$ 的有向曲率为 $\kappa_c=\det(c',c'')/|c'|^3$. 平移不改变两阶导数, 故曲率不变. 对正交矩阵 $A$, 有 $|Ac'|=|c'|$ 和 $\det(Ac',Ac'')=\det(A)\det(c',c'')$. 因而旋转 $A\in SO(2)$ 不改变有向曲率, 反射使其变号; 无向曲率在两者下均不变.

伸缩要求 $r\ne0$, 此时
\[
 \kappa_{rc}=\frac{r^2\det(c',c'')}{|r|^3|c'|^3}
 =\frac{\kappa_c}{|r|}.
\]
若通常约定 $r>0$, 这就是 $\kappa_c/r$. 当 $r<0$ 时, 变换还含一个半周旋转, 所以不能错误地把有向曲率除以负数. $r=0$ 把曲线压成一点, 新曲线不正则, 曲率没有定义.
\end{solution}

\begin{exercise}\label{ps:1-2}
原题PS1.2, 5分. 计算椭圆的曲率, 并确定最大值和最小值的位置.
\end{exercise}
\begin{solution}\label{pssol:1-2}
复用第1章已有的椭圆计算, 并在这里给出原题所需的完整极值结论. 经刚体变换可写 $c(t)=(a\cos t,b\sin t)$, $a\ge b>0$, 沿逆时针定向. 因 $\det(c',c'')=ab$, 得
\[
 \kappa(t)=\frac{ab}{(a^2\sin^2t+b^2\cos^2t)^{3/2}}.
\]
当 $a>b$ 时, 分母在 $t=0,\pi$ 处最小, 故长轴两端 $(\pm a,0)$ 的曲率最大, 为 $a/b^2$; 在 $t=\pi/2,3\pi/2$ 处最大, 故短轴两端 $(0,\pm b)$ 的曲率最小, 为 $b/a^2$. 若 $a=b$, 处处曲率为 $1/a$, 每一点都是最大值和最小值点. 反向参数使有向曲率整体变号; 上述极值大小指按逆时针定向的正曲率, 也即无向曲率.
\end{solution}

\begin{exercise}\label{ps:1-3}
原题PS1.3, 10分. 设 $c$ 是正则曲线, 且对所有参数 $s$ 有 $|c(s)|\le1$. 若有一点 $t$ 满足 $|c(t)|=1$, 证明该点的曲率满足 $|\kappa(t)|\ge1$.
\end{exercise}
\begin{solution}\label{pssol:1-3}
此结论针对参数区间内部的接触点; 若正则曲线的定义采用开区间, 这一条件自动满足. 在该点附近改用弧长参数 $u$, 令接触点为 $u=0$, 写 $\alpha(0)=c(t)$. 函数 $q(u)=|\alpha(u)|^2$ 在0取得局部最大值1, 所以
\[
 q''(0)=2\bigl(|\alpha'(0)|^2+\langle\alpha(0),\alpha''(0)\rangle\bigr)
 =2\bigl(1+\langle\alpha(0),\alpha''(0)\rangle\bigr)\le0.
\]
由 $|\alpha(0)|=1$ 及 Cauchy--Schwarz 不等式, 得 $|\alpha''(0)|\ge1$. 单位速度平面曲线满足 $|\alpha''|=|\kappa|$, 因而得证. 单位圆说明界1可以达到.

若允许接触点只是参数区间端点, 原句需要这一条件批注: $c(s)=(1-s,0)$, $0\le s\le1/2$, 是单位圆盘内的正则直线段, 在 $s=0$ 接触单位圆, 但曲率为0. 端点处没有可用于上述论证的双侧局部最大值二阶条件.
\end{solution}
